% GENERATED FILE. Edit canonical lessons, then run npm run generate:books.
% canonical-source-hash: fnv1a64:4dc153c0fef2c020
% canonical-lessons: ZH-C253-yeye, ZH-C253-nainai, ZH-C253-xiaohai, ZH-W253-hai-child, ZH-C253-zinu, ZH-C253-qinqi, ZH-R253-review-248-253

\chapter{Grandparents and Relatives}
\label{ch:zh-grandparents-and-relatives}
\hlchaptermodality{253}

\begin{chapteropening}
\textbf{By the end of this chapter:} \emph{I can talk about my grandparents, a child, sons and daughters, and relatives.}

Complete the last lesson of chapter 253: I can talk about my grandparents, a child, sons and daughters, and relatives.
\end{chapteropening}

\section[yéye]{yéye — grandpa on the father's side}

\label{lesson:ZH-C253-yeye}



\begin{quote}
Before the new one: say the Mandarin for disappointed, then the Mandarin for satisfied.
\end{quote}

\subsection*{You'll want to know: yéye}
\emph{yéye} — \textquotedblleft{}grandpa on the father's side\textquotedblright{}.

\begin{quote}
\emph{wǒ yéye jīnnián bāshí suì} — My grandpa is eighty this year.
\end{quote}

Its characters come later; for now it is a word for the ear.

\begin{cousinweb}
\emph{yéye} is one syllable, \emph{yé}, said twice, the second time with no tone. No dictionary we use records an older form, so its history is not settled.
\end{cousinweb}

\begin{grammarlens}[title={What you've built}]
One more word for what I have and can do.
\end{grammarlens}

\subsection*{Your turn}
\begin{itemize}
\raggedright
  \item \emph{Say it:} \emph{yéye}
  \item \emph{Say it:} \emph{yéye}, once more
  \item \emph{Say it:} \emph{wǒ yéye jīnnián bāshí suì} — My grandpa is eighty this year.
  \item \emph{Recall:} say the Mandarin for disappointed, then the Mandarin for satisfied, then say \emph{yéye} again
\end{itemize}

\begin{tcolorbox}[breakable,colback=teal!4,colframe=teal!35!black,title={Before you move on}]
What does yéye mean? (\textquotedblleft{}Grandpa on the father's side\textquotedblright{}.) Say it once more.
\end{tcolorbox}
\section[nǎinai]{nǎinai — grandma on the father's side}

\label{lesson:ZH-C253-nainai}



\begin{quote}
Before the new one: say the Mandarin for satisfied, then the Mandarin for grandpa on the father's side.
\end{quote}

\subsection*{You'll want to know: nǎinai}
\emph{nǎinai} — \textquotedblleft{}grandma on the father's side\textquotedblright{}.

\begin{quote}
\emph{nǎinai hé yéye dōu hěn hǎo} — My grandma and grandpa are both well.
\end{quote}

Its characters come later; for now it is a word for the ear.

\begin{cousinweb}
\emph{nǎinai} is built from \emph{nǎi}, said twice, the second time with no tone. It is inherited from Middle Chinese \emph{nreaX} (Guangyun, in Baxter's transcription), which spells the character another way and glosses it 'milk', the \emph{nǎi} of \emph{niúnǎi}; read \emph{nejX}, the same character is glossed as what the people of Chu called their mother. How it came to mean a father's mother is not settled.
\end{cousinweb}

\begin{grammarlens}[title={What you've built}]
One more word for what I have and can do.
\end{grammarlens}

\subsection*{Your turn}
\begin{itemize}
\raggedright
  \item \emph{Say it:} \emph{nǎinai}
  \item \emph{Say it:} \emph{nǎinai}, once more
  \item \emph{Say it:} \emph{nǎinai hé yéye dōu hěn hǎo} — My grandma and grandpa are both well.
  \item \emph{Recall:} say the Mandarin for satisfied, then the Mandarin for grandpa on the father's side, then say \emph{nǎinai} again
\end{itemize}

\begin{tcolorbox}[breakable,colback=teal!4,colframe=teal!35!black,title={Before you move on}]
What does nǎinai mean? (\textquotedblleft{}Grandma on the father's side\textquotedblright{}.) Say it once more.
\end{tcolorbox}
\section[\texorpdfstring{\zh{小孩} (xiǎohái)}{小孩 (xiǎohái)}]{\zh{小孩} (xiǎohái) — child; kid}

\label{lesson:ZH-C253-xiaohai}



\begin{quote}
Before the new one: say the Mandarin for grandpa on the father's side, then the Mandarin for grandma on the father's side.
\end{quote}

\subsection*{You'll want to know: \zh{小孩}}
\textbf{\zh{小孩}} — \emph{xiǎohái} — \textquotedblleft{}child; kid\textquotedblright{}.

\begin{quote}
\textbf{\zh{那个小孩八岁}} \emph{nàge xiǎohái bā suì} — That kid is eight.
\end{quote}

The new character, in pieces: \zh{孩} — nine strokes: \zh{子} on the left, its last stroke rising; on the right a dot, a bar, a stroke that falls to the left and turns up, two strokes falling to the left, and a dot. The next lesson puts it in your hand.

\begin{cousinweb}
\textbf{\zh{小孩}} (\emph{xiǎohái}) is built from \zh{小} 'small' and \zh{孩} 'a child', both inherited from Middle Chinese, \emph{sjewX} and \emph{hoj} (Guangyun, in Baxter's transcription), which glosses the second 'a newborn child'.
\end{cousinweb}

\begin{grammarlens}[title={What you've built}]
One more word for what I have and can do.
\end{grammarlens}

\subsection*{Your turn}
\begin{itemize}
\raggedright
  \item \emph{Say it:} \emph{xiǎohái}
  \item \emph{Say it:} \emph{xiǎohái}, once more
  \item \emph{Say it:} \emph{nàge xiǎohái bā suì} — That kid is eight.
  \item \emph{Recall:} say the Mandarin for grandpa on the father's side, then the Mandarin for grandma on the father's side, then say \emph{xiǎohái} again
\end{itemize}

\begin{tcolorbox}[breakable,colback=teal!4,colframe=teal!35!black,title={Before you move on}]
What does \zh{小孩} mean? (\textquotedblleft{}Child; kid\textquotedblright{}.) Say it once more.
\end{tcolorbox}
\section[\texorpdfstring{\zh{孩} (hái)}{孩 (hái)}]{\zh{孩} — \zh{子}, then a dot, a bar and four strokes}

\label{lesson:ZH-W253-hai-child}



\begin{quote}
The word first, then the character inside it.

\begin{itemize}
\raggedright
  \item \emph{Recall:} say \emph{xiǎohái} — \textbf{R1}, one lesson back, before it decays
\end{itemize}

\begin{quote}
\zh{子} \zh{小}
\end{quote}

\emph{Read it:} each one aloud, with nothing said first — \emph{zǐ}, \emph{xiǎo}
\end{quote}

\subsection*{Script — \zh{子}, then a dot, a bar and four strokes}
\hlblockfigure{\detokenize{figures/ZH-W253-hai-child-filmstrip.pdf}}{How \zh{孩} is written, stroke by stroke}

\begin{quote}
\zh{孩}
\end{quote}

\textbf{Two pieces, left before right.}

\textbf{The left piece, \zh{子}, three strokes.} A bar that turns down to the left, a vertical with a hook, then the last stroke: on the left of a character it rises to the right instead of lying flat.

\textbf{The right piece, six strokes.} A dot on top. A bar under it. A stroke that falls to the left and turns up to the right. A long stroke falling to the left. A shorter stroke falling to the left under it. Then a dot at the lower right.

3 + 6 = nine.

\subsection*{Your turn}
\begin{itemize}
\raggedright
  \item \emph{Write it:} \zh{孩}, all nine, with the model in front of you
  \item \emph{Cover:} write it again from memory, and count the pieces, not the strokes
\end{itemize}

\begin{tcolorbox}[breakable,colback=teal!4,colframe=teal!35!black,title={Before you move on}]
How many strokes in \zh{孩}? (\textbf{Nine}.) Which piece is on the left? (\textbf{\zh{子}}.) What changes in \zh{子} here? (\textbf{Its last stroke rises to the right}.)

Source: \href{https://raw.githubusercontent.com/chanind/hanzi-writer-data/68d10a4b21150cae5e1ebbd223eed289cf32d90c/data/\%E5\%AD\%A9.json}{Hanzi Writer Data \zh{孩}, pinned stroke paths and medians}.
\end{tcolorbox}
\section[\texorpdfstring{\zh{子女} (zǐnǚ)}{子女 (zǐnǚ)}]{\zh{子女} (zǐnǚ) — sons and daughters; children}

\label{lesson:ZH-C253-zinu}



\begin{quote}
Before the new one: say the Mandarin for grandma on the father's side, then the Mandarin for a kid.
\end{quote}

\subsection*{You'll want to know: \zh{子女}}
\textbf{\zh{子女}} — \emph{zǐnǚ} — \textquotedblleft{}sons and daughters; children\textquotedblright{}.

\begin{quote}
\textbf{\zh{他们有三个子女}} \emph{tāmen yǒu sān ge zǐnǚ} — They have three children.
\end{quote}

Two third tones meet in this word: zǐ is said with a rising tone before nǚ, and the spelling keeps both marks.

\begin{cousinweb}
\textbf{\zh{子女}} (\emph{zǐnǚ}) is built from \zh{子} 'a son' and \zh{女} 'a daughter', both inherited from Middle Chinese, \emph{tsiX} and \emph{nrjoX} (Guangyun, in Baxter's transcription).
\end{cousinweb}

\begin{grammarlens}[title={What you've built}]
One more word for what I have and can do.
\end{grammarlens}

\subsection*{Your turn}
\begin{itemize}
\raggedright
  \item \emph{Say it:} \emph{zǐnǚ}
  \item \emph{Say it:} \emph{zǐnǚ}, once more
  \item \emph{Say it:} \emph{tāmen yǒu sān ge zǐnǚ} — They have three children.
  \item \emph{Recall:} say the Mandarin for grandma on the father's side, then the Mandarin for a kid, then say \emph{zǐnǚ} again
\end{itemize}

\begin{tcolorbox}[breakable,colback=teal!4,colframe=teal!35!black,title={Before you move on}]
What does \zh{子女} mean? (\textquotedblleft{}Sons and daughters; children\textquotedblright{}.) Say it once more.
\end{tcolorbox}
\section[qīnqi]{qīnqi — relative; family member}

\label{lesson:ZH-C253-qinqi}



\begin{quote}
Before the new one: say the Mandarin for a kid, then the Mandarin for sons and daughters.
\end{quote}

\subsection*{You'll want to know: qīnqi}
\emph{qīnqi} — \textquotedblleft{}relative; family member\textquotedblright{}.

\begin{quote}
\emph{tā shì wǒ de qīnqi} — He's a relative of mine.
\end{quote}

Its characters come later; for now it is a word for the ear.

\begin{cousinweb}
\emph{qīnqi} is built from \emph{qīn} 'close' and \emph{qi}, both inherited from Middle Chinese, \emph{tshin} and \emph{tshek} (Guangyun, in Baxter's transcription), which glosses the first 'to love' and 'near', and the second with this very word. In the word the second loses its tone.
\end{cousinweb}

\begin{grammarlens}[title={What you've built}]
One more word for what I have and can do.
\end{grammarlens}

\subsection*{Your turn}
\begin{itemize}
\raggedright
  \item \emph{Say it:} \emph{qīnqi}
  \item \emph{Say it:} \emph{qīnqi}, once more
  \item \emph{Say it:} \emph{tā shì wǒ de qīnqi} — He's a relative of mine.
  \item \emph{Recall:} say the Mandarin for a kid, then the Mandarin for sons and daughters, then say \emph{qīnqi} again
\end{itemize}

\begin{tcolorbox}[breakable,colback=teal!4,colframe=teal!35!black,title={Before you move on}]
What does qīnqi mean? (\textquotedblleft{}Relative; family member\textquotedblright{}.) Say it once more.
\end{tcolorbox}
\section[(dialogue)]{Looking back: from A New Term to Grandparents and Relatives}

\label{lesson:ZH-R253-review-248-253}



\begin{quote}
Say the Mandarin for a school term and a year group at school.
\end{quote}

\subsection*{Your turn}
\begin{itemize}
\raggedright
  \item \emph{Say it:} a school term, a year group at school, a head teacher, senior high school, a pupil's parent
  \item \emph{Say it:} to prepare a lesson, to finish school for the day, how a character is read, marks at school, hard-working
  \item \emph{Say it:} a hobby, an interest, a novel, a story, to play chess
  \item \emph{Say it:} measure word for essays, a school essay, to write essays and stories, a new character, to revise a piece of work
  \item \emph{Say it:} in high spirits, to feel, heartbroken, disappointed, satisfied
  \item \emph{Say it:} grandpa on the father's side, grandma on the father's side, a kid, sons and daughters, a relative
  \item \emph{Read it:} \zh{校} \zh{高} \zh{读} \zh{作} \zh{写} \zh{孩} — each one aloud, then a word that holds it
\end{itemize}

\begin{tcolorbox}[breakable,colback=teal!4,colframe=teal!35!black,title={Before you move on}]
A school term? (\textbf{\zh{学期}}, \emph{xuéqī}.) A year group at school? (\emph{niánjí}.) A head teacher? (\textbf{\zh{校长}}, \emph{xiàozhǎng}.) Senior high school? (\textbf{\zh{高中}}, \emph{gāozhōng}.) A pupil's parent? (\textbf{\zh{家长}}, \emph{jiāzhǎng}.) To prepare a lesson? (\emph{yùxí}.) To finish school for the day? (\emph{fàngxué}.) How a character is read? (\textbf{\zh{读音}}, \emph{dúyīn}.) Marks at school? (\emph{chéngjì}.) Hard-working? (\emph{nǔlì}.) A hobby? (\emph{àihào}.) An interest? (\emph{xìngqù}.) A novel? (\textbf{\zh{小说}}, \emph{xiǎoshuō}.) A story? (\emph{gùshi}.) To play chess? (\emph{xiàqí}.) Measure word for essays? (\emph{piān}.) A school essay? (\textbf{\zh{作文}}, \emph{zuòwén}.) To write essays and stories? (\textbf{\zh{写作}}, \emph{xiězuò}.) A new character? (\textbf{\zh{生字}}, \emph{shēngzì}.) To revise a piece of work? (\emph{xiūgǎi}.) In high spirits? (\textbf{\zh{开心}}, \emph{kāixīn}.) To feel? (\emph{gǎndào}.) Heartbroken? (\emph{shāngxīn}.) Disappointed? (\emph{shīwàng}.) Satisfied? (\emph{mǎnyì}.) Grandpa on the father's side? (\emph{yéye}.) Grandma on the father's side? (\emph{nǎinai}.) A kid? (\textbf{\zh{小孩}}, \emph{xiǎohái}.) Sons and daughters? (\textbf{\zh{子女}}, \emph{zǐnǚ}.) A relative? (\emph{qīnqi}.)
\end{tcolorbox}
